\subsection{Preliminaries of Single-Agent Actor and Critic}

RL was born as a tool to optimize decision-making and control, through the experience accumulated by an agent during sequential interactions with an environment, following a trial-and-error strategy. Specifically, at time step $t$, conditionally to the state $s_t \in \mathcal{S}$ of the environment, the agent selects an action $a_t \in \mathcal{A}$ according to its current policy $\pi$ and executes $a_t$ in the environment. At time step $t+1$, based on its reaction to $a_t$, the environment switches to state $s_{t+1}$ and gives a reward $r_t$ back to the agent. When states are partially-observable, an agent can only collect an observation $o_t \in \mathcal{O}$, which contains partial information of the global state $s_t$. Therefore, an action $a_t$ is selected based on the policy $\pi$ and conditionally to the current observation $o_t$. During the interactions, the agent adjusts the policy $\pi$ so as to maximize the long-term cumulative reward, namely the \emph{expected return} $\mathbb{E}_{\pi}[G_t] = \mathbb{E}[\sum^\infty_{k=t+1}\gamma^{k-t-1}r_k]$, where $\gamma$ is a discount factor controlling the importance of a future reward to the current utility.

A first type of learning approaches resort to an action-value function $Q(s_t,a_t)=\mathbb{E}_{\pi}[G_t|s_t, a_t]$ that corresponds to the expected return from state $s_t$, taking action $a_t$ and following policy $\pi$ afterwards.
An RL agent iteratively estimates this action-value function and selects at each step the action with the maximum function value in the current state. This is the fundamental idea of \textit{value-based} RL approaches. $Q^{\psi}(s_t,a_t)$ can be approximated as a $\psi$-parametric function of state and action, which can take the form of a deep neural network (DNN) with weights $\psi$. Parameters $\psi$ can be estimated via temporal-difference approaches by minimizing the regression loss $\mathcal{L}_{Q}(\psi)=\mathbb{E}_{(s_t,a_t,r_t,s_{t+1})\sim H}[(Q^{\psi}(s_t, a_t)-y)^2]$, where $y=r_t +\gamma\mathbb{E}_{a_{t+1}\sim\pi(s_{t+1})}[Q^{\bar{\psi}}(s_{t+1},a_{t+1})]$ is the updated return, $Q^{\bar{\psi}}$ is a moving average of past Q functions, and $H$ is a replay buffer storing past agent-environment interaction data tuples, including states, actions, and rewards.

A second family of learning techniques face the problem from a different perspective, and they are called \textit{policy-gradient} RL approaches. They see a policy as a function indicating the probability of selecting action $a_t$ in state $s_t$, parameterized with vector $\theta$, $\pi_{\theta}(a_t|s_t)=\Pr_{\theta}\{a_t|s_t\}$, which can be represented by a DNN as well, with $\theta$ as connection weights.
Parameters $\theta$ are updated by applying approximate gradient ascent to $\mathbb{E}[G_t]$, thus considering $\nabla_\theta\mathbb{E}[G_t]$, whose unbiased estimate is $\nabla_\theta \log \pi_{\theta}(a_t|s_t)G_t$. Further, $G_t$ can be approximated by its expectation $\mathbb{E}_{\pi}[G_t|s_t, a_t]$, which corresponds to the action-value function $Q^{\psi}(s_t, a_t)$. Finally, to reduce the estimate variance during updates, a state-dependent baseline $b(s_t)$ value is often subtracted from the unbiased estimate, which leads to the gradient $\nabla_\theta \log \pi_{\theta}(a_t|s_t)(Q^{\psi}(s_t, a_t)-b(s_t))$, where $Q^{\psi}(s_t, a_t)-b(s_t)$ is the \textit{advantage} of selecting action $a_t$ over other actions in state $s_t$.

The previous two paragraphs have briefly outlined the two main components of \textit{Actor-Critic} techniques \cite{mnih2016asynchronous}, which have emerged as one of the best-performing RL approaches. Indeed, they are derived from a policy-gradient approach, but incorporate the strengths of a value-based approach. In particular, the \emph{critic} part estimates the action-value function based on past interactions, thus generating $Q^{\psi}(s_t, a_t)$ values, while the \emph{actor} part updates the policy $\pi_{\theta}(a_t|s_t)$ according to the gradient direction, which in turn depends on the action-value function generated by the critic\footnote{Following the convention in Actor-Critic techniques, we will interchangeably refer to, respectively, critic function, action-value function or Q-value function as $Q$ and actor function or policy function as $\pi$ in the remainder of the article.}. Note that when considering partially-observable states, the policy becomes $\pi_{\theta}(a_t|o_t)$, which maps partial observation $o_t$ into a probability distribution over the action set. Similarly, the action-value function becomes $Q^{\psi}(o_t, a_t)$.



\subsection{Complexity and Scalability Analysis}

Let $U_{hid}$ denote the hidden dimension of each layer, then the total number of units in the central Q-value DNN is $U_Q=3U_{hid} + 5U_{hid} \cdot N$, where $N$ is the number of the agents. Specifically, $3U_{hid}$ counts the number of units in the shared attention part that consists of ``key", ``value" and ``agent selection", while $5U_{hid}$ corresponds to two embedding layers and $f^{(i)}$'s two layers where the first contains $2U_{hid}$ units and the second contains $U_{hid}$ units. Then, considering that each agent's policy network consists of 3 layers, the number of units in each agent's policy DNN is $3U_{hid}$, hence the total number of policy units is $U_A = 3U_{hid}\cdot N$. Therefore, the overall number of units is in the order of $\mathcal{O}(U_Q+U_A) = \mathcal{O}(3U_{hid}+8U_{hid}\cdot N)$, which approximates $\mathcal{O}(N)$, as $U_{hid}$ is a constant term potentially taking values of $\{8,16,32,64,128,256,512\}$ (e.g., $U_{hid}=128$ throughout the experiments in this article).
Based on the above analysis, the size of the DNN is only proportional to the number of IAB-nodes, which is usually smaller than 10, as indicated by 3GPP specifications.
From a different perspective, the action space of each agent only depends on the number of sectors of each antenna panel, which is constant in a specific environment, regardless of how many users exist in the system. Hence, the size of the action space, similar to the DNN's size, does not depend on the number of UEs.
In conclusion, the size of the whole training system depends only on the number of IAB-nodes, which is small in practice. Therefore, we can assume our proposed approach to be scalable.

In the policy execution phase, the agents independently apply the individual policies without any need to communicate with the central entity or consider other agents' actions. An agent simply infers its optimal action from its local observation. Therefore, the computational complexity of the proposed MARL-based approach during its execution phase is determined only by the complexity of each local policy.